\documentclass[11pt]{article}

\usepackage[T1]{fontenc}
\usepackage{lmodern}
\usepackage[margin=0.92in]{geometry}
\usepackage{microtype}
\usepackage{amsmath,amssymb,amsthm,mathtools}
\usepackage{booktabs,longtable,array,tabularx}
\usepackage{graphicx}
\usepackage{xcolor}
\usepackage{enumitem}
\usepackage{fancyhdr}
\usepackage{hyperref}
\usepackage[nameinlink,noabbrev]{cleveref}
\usepackage{listings}
\usepackage{caption}
\usepackage{bm}
\usepackage{needspace}

\graphicspath{{figures/}}
\allowdisplaybreaks

\hypersetup{
  pdftitle={Projector-Resolved Seed Calculus II},
  pdfauthor={Komorebi P.},
  pdfsubject={Polyphase refinement, weighted transfer, and dyadic-odd resolution},
  colorlinks=true,
  linkcolor=black,
  citecolor=black,
  urlcolor=black
}

\pagestyle{fancy}
\fancyhf{}
\fancyhead[L]{Projector-Resolved Seed Calculus II}
\fancyhead[R]{Komorebi P.}
\fancyfoot[C]{\thepage}
\setlength{\headheight}{14pt}

\newtheorem{theorem}{Theorem}[section]
\newtheorem{proposition}[theorem]{Proposition}
\newtheorem{lemma}[theorem]{Lemma}
\newtheorem{corollary}[theorem]{Corollary}
\theoremstyle{definition}
\newtheorem{definition}[theorem]{Definition}
\newtheorem{example}[theorem]{Example}
\newtheorem{construction}[theorem]{Construction}
\theoremstyle{remark}
\newtheorem{remark}[theorem]{Remark}
\newtheorem{boundary}[theorem]{Boundary}
\newtheorem{certificate}[theorem]{Numerical Certificate}

\newcommand{\Seed}{\mathfrak S}
\newcommand{\cX}{\mathcal X}
\newcommand{\cA}{\mathcal A}
\newcommand{\cB}{\mathcal B}
\newcommand{\cC}{\mathcal C}
\newcommand{\cD}{\mathcal D}
\newcommand{\cF}{\mathcal F}
\newcommand{\cG}{\mathcal G}
\newcommand{\cO}{\mathcal O}
\newcommand{\cP}{\mathcal P}
\newcommand{\cR}{\mathcal R}
\newcommand{\cS}{\mathcal S}
\newcommand{\cT}{\mathcal T}
\newcommand{\End}{\operatorname{End}}
\newcommand{\Hom}{\operatorname{Hom}}
\newcommand{\im}{\operatorname{im}}
\newcommand{\Ker}{\operatorname{Ker}}
\newcommand{\Spec}{\operatorname{Spec}}
\newcommand{\Res}{\operatorname{Res}}
\newcommand{\Hol}{\operatorname{Hol}}
\newcommand{\Tr}{\operatorname{Tr}}
\newcommand{\diag}{\operatorname{diag}}
\newcommand{\id}{\mathrm{id}}
\newcommand{\bbk}{\mathbb{k}}
\newcommand{\N}{\mathbb N}
\newcommand{\Z}{\mathbb Z}
\newcommand{\C}{\mathbb C}
\newcommand{\Q}{\mathbb Q}
\newcommand{\F}{\mathbb F}
\newcommand{\OO}{\mathcal O}
\newcommand{\Frac}{\operatorname{Frac}}
\newcommand{\val}{\operatorname{val}}
\newcommand{\wt}{\operatorname{wt}}
\newcommand{\lcm}{\operatorname{lcm}}
\newcommand{\packet}[2]{#1_{#2}}
\newcolumntype{L}[1]{>{\raggedright\arraybackslash}p{#1}}
\newcolumntype{Y}{>{\raggedright\arraybackslash}X}

\lstset{
  basicstyle=\ttfamily\footnotesize,
  breaklines=true,
  frame=single,
  showstringspaces=false,
  columns=fullflexible,
  keepspaces=true
}

\title{\textbf{Projector-Resolved Seed Calculus II}\\[0.45em]
\large Polyphase Refinement, Weighted Transfer, and Dyadic--Odd Resolution}
\author{Komorebi P.}
\date{Version 1.2.1 --- September 3, 2026}

\begin{document}
\maketitle

\begin{abstract}
This paper develops the refinement layer of the projector-resolved seed
calculus. A packet decomposition is promoted from one coordinate choice to a
coherent tower of exact coordinate systems: coordinate changes satisfy
identity and cocycle laws, and typed operator matrices are natural under
simultaneous source and target refinement. Coarse packets are sums of fine
descendants, operator blocks aggregate exactly across levels, and residue
selectors admit both division-free coefficient formulas and finite-Fourier
realizations over a commutative semiring from a finite orthogonality table.
For $N=2^aM$ with $M$ odd, the Chinese remainder theorem separates dyadic and
odd packet addresses. Weighted transfer inversion is proved in a certified
finite quotient and in a complete normed ring under a strict contraction;
the finite theorem constructs the signed, cumulative-weight-pruned word sum
and permits a sharper length cutoff from a uniform lower edge weight. For a
finite positive-weight multigraph, Lean proves the first-edge recurrence for
the enumerated walk series and identifies it with the automatically generated
determinant/adjugate expression, yielding rationality and a recurrence without
a supplied rationality certificate. Minimal evaluated B\'ezout hypotheses
give a quasi-idempotent numerator and an idempotent separator, while explicit
lattice-unit data gives an upward-closed separation-weight bound. Concrete
$\mathbb Z/p^L\mathbb Z$ Smith fibres have the expected divisibility criterion
and exact cardinality; the existence of Smith normal form itself remains
classical input. The rational sign split has exact $2$-local denominator
weight one, equal to the $2$-adic valuation of its computed resultant, and
exact local weight zero at odd primes. Typed linear defects propagate through
powers, polynomial actions, and explicitly witnessed resolvents, with proved
operator-norm majorants. The companion release contains thirteen generated
figures, a versioned Python implementation with forty-eight tests, and a Lean
4/Mathlib formalization of the frozen twenty-five-item contract. These checks
establish the stated formal contract, not novelty, priority, or external peer
review.
\end{abstract}

\begin{center}
\fbox{\begin{minipage}{0.92\textwidth}
\textbf{Scope.} Paper I \cite{komorebi2026projectorI} fixed the typed seed, operator, relation, packet,
completion, and reconstruction language. Paper II assumes that foundation and
studies exact refinement, coordinate-change coherence, coefficient and phase
selectors, certified weighted inverses, finite weighted-walk enumeration and
rationality, evaluated B\'ezout separators, Smith-coordinate fibres, and typed
linear defects. The transfer-zeta construction is automatic only for the
finite, positive-integral-weight graph model stated here. Universal
complete-filtered-algebra convergence, construction of a general polynomial
resultant certificate, and the structure theorem producing Smith normal form
remain cited or explicit input rather than being hidden inside stronger
claims. The concrete sign-factor resultant is computed separately. Nonlinear
Fa\`a di Bruno trees and general obstruction packets remain outside this paper.
\end{minipage}}
\end{center}

\tableofcontents
\listoffigures
\listoftables
\clearpage

\section{The refinement problem}

Paper I \cite{komorebi2026projectorI} resolved a generated object $x\in X$ by one complete idempotent family
\[
P_\alpha P_\beta=\delta_{\alpha\beta}P_\alpha,
\qquad
\sum_{\alpha\in I}P_\alpha=I_X,
\]
with exact reconstruction
\[
x=\sum_{\alpha\in I}P_\alpha x.
\]
A single decomposition, however, does not yet express scale. The present problem is to place these packet systems into a coherent tower
\[
\cP_0\prec\cP_1\prec\cP_2\prec\cdots
\]
so that every coarse packet is the exact sum of its fine descendants and every operator can be transported between levels without changing the underlying object.

The production language now becomes
\[
\boxed{
\text{seed}
\to
\text{address}
\to
\text{operate}
\to
\text{project}
\to
\text{refine}
\to
\text{weight}
\to
\text{complete}
\to
\text{reconstruct}.
}
\]

Three questions organize the paper.
\begin{enumerate}[label=(\roman*)]
\item When is a fine packet family an exact refinement of a coarse family?
\item When can an operator inverse be reconstructed from weighted packet-transfer words?
\item Why does the dyadic sector sometimes require additional depth even when the corresponding odd-prime packets separate immediately?
\end{enumerate}

\begin{figure}[htbp]
\centering
\includegraphics[width=0.94\textwidth]{01_nested_dyadic_refinement.pdf}
\caption{A nested residue hierarchy. Each residue modulo $2$ is the union of two residues modulo $4$, and each residue modulo $4$ is the union of two residues modulo $8$.}
\label{fig:refinement-tree}
\end{figure}

\subsection{Status of the statements}

The paper uses the following labels.
\begin{itemize}
\item \textbf{Theorem/Proposition}: proved under the stated hypotheses.
\item \textbf{Classical input}: an established algebraic fact used without claiming priority.
\item \textbf{Worked example}: an exact realization of the definitions.
\item \textbf{Numerical certificate}: a finite computational check of the implementation.
\item \textbf{Boundary}: a claim intentionally not made.
\end{itemize}

The release freezes a theorem-completeness denominator of twenty-five items:
the twenty-one numbered definitions, propositions, theorems, and corollaries
in the main text; the exact Smith-phase identity in
\cref{sec:smith-realization}; and the three algorithms in
\cref{sec:algorithms}. Each item is mapped to one or more compiled Lean
declarations. ``Lean-complete'' below always means complete relative to this
revised, scope-explicit contract. Some rows still take application-specific
certificates, and their boundaries are stated individually. It is not a claim that Lean proves the
historical motivation, examples, floating-point residuals, or bibliographic
priority. The formalization uses Lean 4 and Mathlib
\cite{demoura2015lean,mathlib2020}.

\section{Nested packet systems}

\subsection{Refinement data}

Let $X$ be a module over a commutative ring $R$ or a topological vector space over a field. At level $\ell$, let
\[
\cP_\ell=\{P_{\ell,\alpha}\}_{\alpha\in I_\ell}
\]
be a finite complete family of pairwise orthogonal idempotents.

\begin{definition}[Refinement map]\label{def:refinement-map}
A refinement map from level $\ell+1$ to level $\ell$ is a surjection
\[
\rho_{\ell+1,\ell}:I_{\ell+1}\longrightarrow I_\ell
\]
such that
\[
\boxed{
P_{\ell,\alpha}
=
\sum_{\rho_{\ell+1,\ell}(\beta)=\alpha}
P_{\ell+1,\beta}
}
\]
for every $\alpha\in I_\ell$.
\end{definition}

The relation $\beta\succ\alpha$ means that $\beta$ is a fine child of $\alpha$.

\begin{theorem}[Nested refinement theorem]\label{thm:nested-refinement}
Suppose $\cP_{\ell+1}$ refines $\cP_\ell$. Then:
\begin{enumerate}[label=(\alph*)]
\item every fine packet lies inside its parent:
\[
P_{\ell,\alpha}P_{\ell+1,\beta}
=
P_{\ell+1,\beta}P_{\ell,\alpha}
=
\begin{cases}
P_{\ell+1,\beta},&\beta\succ\alpha,\\
0,&\text{otherwise};
\end{cases}
\]
\item every coarse packet vector is reconstructed uniquely from its fine children;
\item refinement is transitive: if $\cP_{\ell+2}$ refines $\cP_{\ell+1}$ and $\cP_{\ell+1}$ refines $\cP_\ell$, then $\cP_{\ell+2}$ refines $\cP_\ell$.
\end{enumerate}
\end{theorem}

\begin{proof}
Multiply
\[
P_{\ell,\alpha}
=
\sum_{\gamma\succ\alpha}P_{\ell+1,\gamma}
\]
by $P_{\ell+1,\beta}$. Orthogonality at the fine level leaves exactly the term $P_{\ell+1,\beta}$ when $\beta\succ\alpha$, and zero otherwise. The reconstruction statement follows by applying the identity to $x$. For transitivity, substitute the child decomposition of each $P_{\ell+1,\beta}$ into the parent decomposition of $P_{\ell,\alpha}$ and collect all grandchildren whose composite parent is $\alpha$.
\end{proof}

\subsection{Analysis, synthesis, refinement, and coarsening}

Define the packet-coordinate space
\[
X_\ell^{\mathrm{pkt}}
=
\bigoplus_{\alpha\in I_\ell}P_{\ell,\alpha}X.
\]
The analysis and synthesis maps are
\[
\cA_\ell x=(P_{\ell,\alpha}x)_\alpha,
\qquad
\cS_\ell(x_\alpha)_\alpha=\sum_\alpha x_\alpha.
\]
They satisfy
\[
\cS_\ell\cA_\ell=I_X,
\qquad
\cA_\ell\cS_\ell=I_{X_\ell^{\mathrm{pkt}}}.
\]

\begin{definition}[Exact coordinate refinement]\label{def:coordinate-refinement}
Define
\[
J_{\ell+1,\ell}=\cA_{\ell+1}\cS_\ell,
\qquad
C_{\ell,\ell+1}=\cA_\ell\cS_{\ell+1}.
\]
The map $J$ resolves coarse packet coordinates into fine coordinates, while $C$ regroups fine coordinates by parent.
\end{definition}

\begin{proposition}[Perfect reconstruction across levels]\label{prop:perfect-reconstruction}
The maps $J_{\ell+1,\ell}$ and $C_{\ell,\ell+1}$ are mutual inverses:
\[
C_{\ell,\ell+1}J_{\ell+1,\ell}=I,
\qquad
J_{\ell+1,\ell}C_{\ell,\ell+1}=I.
\]
\end{proposition}

\begin{proof}
Use $\cS_m\cA_m=I_X$ at the intermediate level:
\[
C J
=\cA_\ell\cS_{\ell+1}\cA_{\ell+1}\cS_\ell
=\cA_\ell\cS_\ell=I.
\]
The other identity is identical with the levels exchanged.
\end{proof}

These coordinate changes are coherent through every finite tower. If three
complete packet systems are denoted by $a,b,c$, then
\[
J_{c,b}J_{b,a}=J_{c,a},
\qquad J_{a,a}=I,
\qquad J_{b,a}^{-1}=J_{a,b}.
\]
These laws hold for arbitrary exact packet systems representing the same
ambient module; the parent maps are needed for the componentwise
parent--child formulas, not for coordinate coherence itself.

Moreover, let $P_a,P_b$ be packet systems on $X$, let $Q_a,Q_b$ be packet
systems on $Y$, and let $T:X\to Y$ be a typed linear map. Its exact coordinate
representation
\[
\widehat T_{Q,P}=\cA_QT\cS_P
\]
preserves identities and typed composition, and the naturality square is
\[
\boxed{
J^Y_{Q_b,Q_a}\,\widehat T_{Q_a,P_a}
=
\widehat T_{Q_b,P_b}\,J^X_{P_b,P_a}.
}
\]
The identity, cocycle, composition, inverse, and naturality laws are all proved
in Lean. The separate refinement hypotheses supply the local parent--child
component formulas.

This result should not be confused with dimensional reduction. Both packet-coordinate spaces represent the full object. Refinement changes coordinates; it does not discard fine information.

\subsection{Operator blocks under refinement}

For $T\in\End(X)$, define
\[
T^{(\ell)}_{\beta\alpha}
=P_{\ell,\beta}TP_{\ell,\alpha}.
\]

\begin{theorem}[Transfer-block refinement]\label{thm:transfer-block-refinement}
If $\cP_{\ell+1}$ refines $\cP_\ell$, then
\[
\boxed{
T^{(\ell)}_{ba}
=
\sum_{\substack{\beta\succ b\\\alpha\succ a}}
T^{(\ell+1)}_{\beta\alpha}.
}
\]
Consequently the complete operator is reconstructed at every level:
\[
T=\sum_{a,b}T^{(\ell)}_{ba}
=\sum_{\alpha,\beta}T^{(\ell+1)}_{\beta\alpha}.
\]
\end{theorem}

\begin{proof}
Expand both coarse projectors into their children:
\[
P_{\ell,b}TP_{\ell,a}
=
\left(\sum_{\beta\succ b}P_{\ell+1,\beta}\right)
T
\left(\sum_{\alpha\succ a}P_{\ell+1,\alpha}\right).
\]
Distribute the finite sums.
\end{proof}

\begin{figure}[htbp]
\centering
\includegraphics[width=0.92\textwidth]{03_refinement_incidence.pdf}
\caption{The incidence matrix $M_{r,s}=1$ when the fine residue $s$ belongs to the coarse residue $r$. Its rows implement exact fine-to-coarse packet aggregation.}
\label{fig:incidence}
\end{figure}

\section{Residue and polyphase selectors}

\subsection{Coefficient selectors over an arbitrary ring}

Let
\[
F(z)=\sum_{n\ge0}a_nz^n
\]
be a formal series over a ring $R$. For $N\ge1$ and $0\le j<N$, define
\[
\boxed{
(P_{N,j}F)(z)
=
\sum_{n\equiv j\,(\mathrm{mod}\,N)}a_nz^n.
}
\]
No division and no roots of unity are required.

\begin{proposition}[Residue projector algebra]\label{prop:residue-algebra}
The coefficient selectors satisfy
\[
P_{N,i}P_{N,j}=\delta_{ij}P_{N,j},
\qquad
\sum_{j=0}^{N-1}P_{N,j}=I.
\]
\end{proposition}

\begin{proof}
Every coefficient address $n$ belongs to exactly one residue class modulo $N$. Applying two selectors preserves the coefficient only when both residue labels agree with $n$.
\end{proof}

\begin{figure}[htbp]
\centering
\includegraphics[width=0.95\textwidth]{02_crt_dyadic_odd_grid.pdf}
\caption{A composite residue packet is addressed by its dyadic and odd coordinates. The displayed mod-$24$ system is the product of mod $8$ and mod $3$.}
\label{fig:crt-grid}
\end{figure}

\subsection{Finite Fourier realization}

Let $U$ act on coefficient sequences by
\[
(UF)(z)=F(\omega z).
\]
If $\omega^N=1$, then $U^N=I$.

\begin{theorem}[Root-of-unity realization]\label{thm:root-of-unity}
Let $R$ be a commutative semiring, let $N>0$, and suppose
$\omega,\omega^\vee,\eta\in R$ satisfy $\omega^N=1$ and the finite
$N\times N$ Fourier orthogonality table
\[
\eta\sum_{k=0}^{N-1}(\omega^\vee)^{jk}\omega^{rk}
=
\begin{cases}
1,&r=j,\\
0,&r\ne j,
\end{cases}
\qquad 0\le j,r<N.
\]
Then
\[
\boxed{
P_{N,j}
=
\eta\sum_{k=0}^{N-1}(\omega^\vee)^{jk}U^k.
}
\]
In particular, if $R$ is a field, $N>0$, and $\omega$ is a primitive
$N$th root of unity, the certificate holds with
$\omega^\vee=\omega^{-1}$ and $\eta=N^{-1}$.
\end{theorem}

\begin{proof}
On the monomial $z^n$, the right-hand side acts by
\[
\eta\sum_{k=0}^{N-1}(\omega^\vee)^{jk}\omega^{nk}.
\]
The identity $\omega^N=1$ reduces this expression to the table entry indexed
by $r=n\bmod N$. The finite certificate therefore gives exactly the
coefficient mask defining $P_{N,j}$.
Over a field, the primitive-root geometric sum supplies the certificate.
\end{proof}

The field or orthogonality condition is essential. A primitive element of
order $N$ in a ring with zero divisors does not by itself force the off-diagonal
geometric sums to vanish. For example, in $\mathbb F_3\times\mathbb F_3$ the
element $(-1,1)$ has order two, but the second component of its nontrivial
character sum is nonzero. The division-free coefficient selector remains valid
without any Fourier certificate. This distinction becomes decisive at the
dyadic prime.

\subsection{Nested residue refinement}

\begin{theorem}[Polyphase refinement]\label{thm:polyphase-refinement}
If $M=Nq$ with $N,q>0$, then
\[
\boxed{
P_{N,j}
=
\sum_{\substack{0\le s<M\\s\equiv j\,(\mathrm{mod}\,N)}}
P_{M,s}.
}
\]
\end{theorem}

\begin{proof}
An address $n$ satisfies $n\equiv j\pmod N$ exactly when its unique residue $s=n\pmod M$ satisfies $s\equiv j\pmod N$.
\end{proof}

Thus the familiar $2\to4\to8$ hierarchy is exact, not approximate:
\[
P_{2,0}=P_{4,0}+P_{4,2}
=P_{8,0}+P_{8,2}+P_{8,4}+P_{8,6},
\]
\[
P_{2,1}=P_{4,1}+P_{4,3}
=P_{8,1}+P_{8,3}+P_{8,5}+P_{8,7}.
\]

\section{Dyadic--odd factorization}

Let
\[
N=2^aM,
\qquad M\ \text{odd}.
\]
The two moduli are coprime.

\begin{theorem}[Dyadic--odd packet factorization]\label{thm:dyadic-odd}
For every residue $j\pmod N$, let
\[
j_2=j\pmod{2^a},
\qquad
j_o=j\pmod M.
\]
Then
\[
\boxed{
P_{N,j}=P_{2^a,j_2}P_{M,j_o}=P_{M,j_o}P_{2^a,j_2}.
}
\]
Conversely, every pair $(j_2,j_o)$ labels one and only one packet modulo $N$.
\end{theorem}

\begin{proof}
A coefficient address $n$ survives the product on the right exactly when
\[
n\equiv j_2\pmod{2^a},
\qquad
n\equiv j_o\pmod M.
\]
The Chinese remainder theorem gives one residue $j\pmod{2^aM}$ satisfying both conditions.
\end{proof}

This theorem justifies separating the coordinate system into
\[
\boxed{
\text{dyadic phase address}
\times
\text{odd phase address}.
}
\]
It does not assert that their arithmetic lifting theories are identical. The coordinate factorization is exact, while the integral projector weights can still differ sharply.

\begin{example}[Mod $12$]
Since $12=4\cdot3$, every residue is identified by
\[
(j\bmod4,j\bmod3).
\]
For instance, the pair $(3,1)$ corresponds to $j=7\pmod{12}$, so
\[
P_{12,7}=P_{4,3}P_{3,1}.
\]
\end{example}

\section{Weighted transfer words}

\subsection{Filtered operator algebras}

Let $\OO$ be a complete discrete valuation ring with uniformizer $\pi$, fraction field $K$, and filtration
\[
F^m\cB=\pi^m\cB
\]
on an $\OO$-algebra $\cB$. Suppose
\[
T=D+H,
\]
where $D$ is invertible and the perturbation $H$ raises filtration. The symbol
$K$ is reserved for the fraction field.

A packet decomposition resolves $H=\sum_{e\in E}H_e$ and assigns a
certified filtration gain
\[
D^{-1}H_e\in F^{w(e)}\cB,
\qquad w(e)\in\Z_{>0}.
\]
Each $e$ is a directed transfer edge. For an ordered transfer word
\[
\gamma=(e_1,\ldots,e_r),
\]
define
\[
W(\gamma)=\sum_{j=1}^rw(e_j),
\qquad
A_\gamma=A_{e_1}\cdots A_{e_r},
\qquad A_e=D^{-1}H_e.
\]

\begin{figure}[htbp]
\centering
\includegraphics[width=0.87\textwidth]{04_weighted_transfer_graph.pdf}
\caption{A finite packet-transfer graph. Edge weights record filtration gain rather than merely the number of operator compositions.}
\label{fig:weighted-graph}
\end{figure}

\begin{theorem}[Certified weighted inverse]\label{thm:weighted-inverse}
Let $R$ be a possibly noncommutative ring, let $D^{-1}$ be a supplied
two-sided inverse of $D$, and put $A=D^{-1}H$.
\begin{enumerate}[label=(\alph*)]
\item If $A^r=0$, then $D+H$ has the two-sided inverse
\[
\boxed{(D+H)^{-1}=\left(\sum_{k=0}^{r-1}(-A)^k\right)D^{-1}.}
\]
\item If $R$ is a complete normed ring and $\lVert A\rVert<1$, then
$\sum_{k\ge0}(-A)^k$ converges and
\[
\boxed{(D+H)^{-1}=\left(\sum_{k\ge0}(-A)^k\right)D^{-1}.}
\]
\end{enumerate}
\end{theorem}

\begin{proof}
Factor
\[
T=D(I+A),
\qquad A=D^{-1}H.
\]
For a finite cutoff $r$, multiplication on either side gives the
noncommutative telescoping identity
\[
(I+A)\sum_{k=0}^{r-1}(-A)^k
=
\left(\sum_{k=0}^{r-1}(-A)^k\right)(I+A)
=1-(-A)^r.
\]
This proves (a). In (b), the strict norm bound makes the geometric series
summable. Continuity of multiplication passes the same identities to the
limit. Multiplication by $D^{-1}$ gives the displayed inverse in both orders.
\end{proof}

\begin{corollary}[Finite precision]\label{cor:finite-precision}
Let $F^\bullet R$ be a multiplicative decreasing filtration and let
$q_L:R\to R_L$ be a ring homomorphism that kills $F^L R$. Suppose
\[
H=\sum_{e\in E}H_e,
\qquad D^{-1}H_e\in F^{w(e)}R,
\qquad w(e)\ge1,
\]
with $E$ finite. Then $q_L(D^{-1}H)^L=0$, the finite sum in
\cref{thm:weighted-inverse}(a) is a two-sided inverse after applying $q_L$,
and its word expansion may discard every word of weight at least $L$.
More explicitly, for a length cutoff $N$, define
\[
S^{\mathrm{pr}}_{L,N}
=
\sum_{n=0}^{N-1}(-1)^n
\sum_{\substack{\gamma\in E^n\\W(\gamma)<L}}A_\gamma.
\]
Lean proves directly that $q_L(S^{\mathrm{pr}}_{L,N})$ equals the mapped
geometric inverse whenever $q_L(A^N)=0$, and hence is a two-sided inverse of
$q_L(I+A)$, where $A=\sum_eA_e=D^{-1}H$. With the supplied two-sided
inverse $D^{-1}$ of $D$, this gives the explicit two-sided inverse
\[
\boxed{q_L\!\left(S^{\mathrm{pr}}_{L,N}D^{-1}\right)}
\quad\text{of}\quad q_L(D+H).
\]
If every edge weight is at least $m$ and $L\le mN$, then the nilpotence
condition follows from the filtration data because
$m|\gamma|\le W(\gamma)$. For $m>0$, the executable choice
$N=\lceil L/m\rceil$ satisfies $L\le mN$; this elementary ceiling inequality
is also proved in Lean. The choice $m=1,N=L$ gives the standard cutoff.
In the DVR example $q_L:R\to R/\pi^L R$, this reads
\[
q_L(T)^{-1}
=
q_L\!\left(\sum_{W(\gamma)<L}
(-1)^{|\gamma|}A_\gamma D^{-1}\right).
\]
\end{corollary}

Here $A_\gamma=(D^{-1}H_{e_1})\cdots(D^{-1}H_{e_r})$ in the written
order. Multiplicativity places it in $F^{W(\gamma)}R$, so the quotient kills
it when $W(\gamma)\ge L$. Positive weights also imply
$|\gamma|\le W(\gamma)$, making the retained search finite.

No unrestricted convergence assertion for every complete filtered algebra is
used in the formal contract. Such an assertion needs a separately specified
topology and completeness theorem. The release proves the finite-quotient
result above and the complete-normed contraction specialization.

\begin{figure}[htbp]
\centering
\includegraphics[width=0.84\textwidth]{06_filtered_inverse_depth.pdf}
\caption{At precision $L$, weighted inversion requires only transfer words whose cumulative weight is below $L$.}
\label{fig:inverse-depth}
\end{figure}

\section{Transfer-zeta functions}

Let $G$ be a finite directed graph with positive integer edge weights. Define the polynomial adjacency matrix
\[
A(t)_{ij}
=
\sum_{e:i\to j}t^{w(e)},
\]
where parallel edges are distinct elements of the finite edge type and are
therefore counted with their full multiplicity.

For fixed start and target vertices $s,t$, define
\[
Z_{s,t}(t)
=
\sum_{n\ge0}a_n(s,t)t^n,
\]
where $a_n(s,t)$ counts directed walks from $s$ to $t$ with total weight $n$.
Vertices and edges may repeat. Thus ``path zeta'' is retained as the name of
the series and Lean API, but no simple-path restriction is intended.

\begin{theorem}[Automatic finite-graph transfer-zeta rationality]\label{thm:transfer-zeta}
Assume the vertex and edge sets are finite and every edge has positive integer
weight. For every $n$, exact enumeration of edge tuples of lengths at most
$n$, followed by endpoint and total-weight filtering, returns $a_n(s,t)$.
These coefficients obey the first-edge recurrence
\[
a_n(s,t)
=
\mathbf 1_{s=t,\,n=0}
+\sum_{\substack{e:\,\operatorname{source}(e)=s\\w(e)\le n}}
a_{n-w(e)}(\operatorname{target}(e),t).
\]
Let
\[
T(t)=I-A(t),\qquad
Q(t)=\det T(t),\qquad
P_{s,t}(t)=\operatorname{adj}(T(t))_{s,t}.
\]
Then $Q(0)=1$, hence $Q\ne0$ and $Q$ is a unit in $\Z[[t]]$, and the
independently enumerated walk series satisfies
\[
\boxed{
Z_{s,t}(t)=Q(t)^{-1}P_{s,t}(t),
\qquad
Q(t)Z_{s,t}(t)=P_{s,t}(t)
}
\quad\text{in }\Z[[t]].
\]
\end{theorem}

\begin{proof}
Every edge has weight at least one, so a word of weight $n$ has length at
most $n$. The enumerator therefore searches a finite set containing every
contributing walk, and its filters are exactly the defining predicates for
$a_n$. Removing the first edge gives the displayed disjoint decomposition
into the empty diagonal walk and the first-edge tails. In matrix form this is
\[
Z=I+AZ,
\qquad (I-A)Z=I.
\]
Positive weights give $A(0)=0$, so $T(0)=I$ and
$Q(0)=\det(I)=1$. The two adjugate identities
\[
T\operatorname{adj}(T)=Q I,
\qquad
\operatorname{adj}(T)T=Q I
\]
remain valid after inclusion into $\Z[[t]]$. Multiplying the adjugate by the
formal-series inverse of $Q$ constructs a two-sided inverse of $T$.
Uniqueness of the inverse identifies that construction with the independently
enumerated matrix $Z$, yielding the entrywise formulas.
\end{proof}

This is the finite transfer-matrix method of \cite{stanley2012}, but here its
determinant, adjugate, normalization, first-edge recurrence, and identification
with the enumerated series are all formalized. The older API accepting an
independently supplied normalized coefficient identity $QZ=P$ is retained for
external certificates; it is no longer needed to obtain rationality for the
finite graph defined above. The Python companion performs an exact symbolic
construction as an independent executable cross-check.

\begin{corollary}[Linear recurrence]\label{cor:zeta-recurrence}
With the canonical polynomials from \cref{thm:transfer-zeta}, write
\[
Q(t)=\sum_kq_kt^k,
\qquad
Z_{s,t}(t)=\sum_{n\ge0}a_n(s,t)t^n.
\]
Then the walk counts satisfy the constant-coefficient recurrence determined by
$Q$ after the finite initial range contributed by $P_{s,t}$:
\[
\boxed{
\sum_{k=0}^{n}q_k a_{n-k}(s,t)=0
\qquad\bigl(n>\deg P_{s,t}\bigr).
}
\]
\end{corollary}

\begin{figure}[htbp]
\centering
\includegraphics[width=0.88\textwidth]{05_path_zeta_coefficients.pdf}
\caption{The walk counts for one weighted graph. The exact generating function is rational, so these counts satisfy a finite linear recurrence.}
\label{fig:zeta-counts}
\end{figure}

\begin{example}[A computed transfer zeta]
For the graph used in the numerical companion, the walk series from packet $0$ to packet $2$ is
\[
Z_{0,2}(t)
=
\frac{t^3}{1-t^2-t^3-2t^4+t^5+t^6+t^7-t^9}.
\]
The first nonzero coefficients are
\[
1,1,1,3,1,5,4,10,7,19,17,37,\ldots
\]
at their corresponding weights. Direct enumeration through weight $30$ agrees coefficient by coefficient with the rational expansion.
\end{example}

\section{Integral spectral packets}

\subsection{B\'ezout projectors}

Let $\OO$ be a DVR, $K=\Frac(\OO)$, and $A\in\End_\OO(\Lambda)$ for a finite free lattice $\Lambda$. Suppose the characteristic polynomial factors over $\OO$ as
\[
\chi_A(x)=f(x)g(x),
\]
where $f$ and $g$ are coprime over $K[x]$.

\begin{theorem}[Evaluated B\'ezout spectral projector]\label{thm:bezout-projector}
Let $F,G,U,V\in\End_K(\Lambda\otimes_\OO K)$ and let $r\in K^\times$.
Suppose
\[
UF+VG=rI,
\qquad FG=0,
\qquad FG=GF,
\qquad GU=UG.
\]
Put $N=VG$. Then
\[
\boxed{N^2=rN.}
\]
Consequently the operator
\[
\boxed{
E=r^{-1}VG
}
\]
is idempotent, acts as the identity on $\ker F$, and acts as zero on
$\ker G$. If $V$ and $G$ preserve $\Lambda$, then
\[
rE=VG\in\End_\OO(\Lambda).
\]
\end{theorem}

\begin{proof}
The B\'ezout identity gives $N=rI-UF$. The two stated commutation relations
and $FG=0$ give
\[
N(UF)=VGUF=VUGF=VUFG=0.
\]
Therefore
\[
N^2=N(rI-UF)=rN.
\]
Since $r\ne0$, $E^2=r^{-2}N^2=E$. If $Fx=0$, the evaluated B\'ezout
identity gives $VGx=rx$, hence $Ex=x$. If $Gx=0$, then $Ex=0$.
The lattice assertion follows directly from $rE=VG$.
\end{proof}

The denominator-free equality $N^2=rN$ uses neither $r\ne0$ nor the other
pairwise commutation relations that hold in the polynomial application. The
normalization to $E$ uses $r\ne0$; action on $\ker F$ uses only the B\'ezout
identity and $r\ne0$, while vanishing on $\ker G$ is immediate from $N=VG$.
The bundled Lean certificate records all pairwise commutations as convenient
evaluated polynomial data, but the minimal theorems expose only the hypotheses
they consume.

For the polynomial application, one supplies $F=f(A)$, $G=g(A)$,
$U=u(A)$, and $V=v(A)$. A B\'ezout or resultant identity
$uf+vg=r$ evaluates to the first certificate, while
$(fg)(A)=0$ supplies the second; when $fg=\chi_A$, this last equality is
Cayley--Hamilton. Standard B\'ezout, resultant, and DVR facts are recalled
from \cite{lang2002,serre1979}. The formal theorem begins with the evaluated
certificate. It does not infer exact action merely from nilpotence of
$f(A)$ or $g(A)$, which would be insufficient.

\begin{definition}[Separation weight]\label{def:separation-weight}
For a rational projector $E$ and a specified $\OO$-lattice $\Lambda$, provided
at least one clearing exponent exists, define
\[
s_{\pi,\Lambda}(E)
=
\min\{s\ge0:\pi^sE(\Lambda)\subseteq\Lambda\}.
\]
If the evaluated certificate arises from integral $v,g$ and
$r=u\pi^s$ with $u^{-1}$ preserving the lattice, then
\cref{thm:bezout-projector} gives the upper bound
\[
s_{\pi,\Lambda}(E)\le s=\val_\pi(r).
\]
Indeed,
\[
\boxed{\pi^sE=u^{-1}VG}
\]
preserves $\Lambda$. The minimal Lean theorem needs only the lattice action of
$u^{-1}$; its bundled lattice-unit wrapper records both $u$ and $u^{-1}$,
rather than treating every nonzero fraction-field scalar as an integral unit.
Existence of a clearing exponent is an explicit hypothesis, not a consequence
of the notation. Once one exists, the least one is obtained by well-ordering.
Clearing exponents are upward closed provided multiplication by $\pi$
preserves the lattice.
\end{definition}

\subsection{Precision delay}

\begin{theorem}[Projector precision delay]\label{thm:precision-delay}
Let $E$ be a rational projector with $\pi^sE(\Lambda)\subseteq\Lambda$. If
\[
x\equiv y\pmod{\pi^{m+s}\Lambda},
\]
then
\[
Ex\equiv Ey\pmod{\pi^m\Lambda}.
\]
Thus $m+s$ input digits suffice to determine $m$ output digits of the projected packet.
\end{theorem}

\begin{proof}
Write $x-y=\pi^{m+s}z$. Since $\pi^sE(z)\in\Lambda$,
\[
E(x-y)=\pi^m\bigl(\pi^sE(z)\bigr)\in\pi^m\Lambda.
\]
\end{proof}

The theorem makes delayed separation quantitative. A packet with weight $s$ is not absent; it requires $s$ additional digits before it can be reconstructed to the same output precision.

\section{The exceptional dyadic sign split}

Consider the integral involution
\[
A=
\begin{pmatrix}
1&1\\
0&-1
\end{pmatrix},
\qquad
A^2=I.
\]
Over a field of characteristic different from two, its sign projectors are
\[
E_+=\frac{I+A}{2}
=
\begin{pmatrix}
1&1/2\\
0&0
\end{pmatrix},
\]
\[
E_-=\frac{I-A}{2}
=
\begin{pmatrix}
0&-1/2\\
0&1
\end{pmatrix}.
\]
They are orthogonal idempotents and sum to $I$.

\begin{figure}[htbp]
\centering
\includegraphics[width=0.88\textwidth]{07_dyadic_delayed_separation.pdf}
\caption{The rational sign projectors for an integral involution. Their denominator has $2$-adic weight one.}
\label{fig:dyadic-projectors}
\end{figure}

\begin{proposition}[Odd-adic versus dyadic behavior]\label{prop:dyadic-sign}
For every prime $p\ne2$, every entry of $E_+$ and $E_-$ is $p$-integral,
and each projector has exact local denominator weight zero. At $p=2$, each
projector has exact local denominator weight one; multiplication by $2$ makes
it integral, while the unscaled matrix is not integral:
\[
s_{2,\Z^2}(E_+)=s_{2,\Z^2}(E_-)=1.
\]
The characteristic factors $x-1$ and $x+1$ have resultant $2$. Its $2$-adic
valuation is one, exactly the two local projector weights above, and its
$p$-adic valuation is zero at every prime $p\ne2$.
\end{proposition}

\begin{proof}
The reduced denominators of the displayed entries are $1$ or $2$, which proves
the local integrality and minimality statements directly. The resultant is
\[
\Res(x-1,x+1)=2.
\]
Thus $v_2(2)=1$ and $v_p(2)=0$ for primes $p\ne2$.
\end{proof}

The two characteristic factors $x-1$ and $x+1$ have the same reduction in
$\F_2[x]$. The rational matrices $E_+$ and $E_-$ themselves cannot be reduced
modulo $2$, because their entries contain $1/2$. Consequently the integral
special fibre does not inherit their separate idempotent splitting, even
though the two projectors remain distinct over $\Q_2$. This is the simplest
model of delayed dyadic separation.

\begin{figure}[htbp]
\centering
\includegraphics[width=0.84\textwidth]{08_odd_prime_separation.pdf}
\caption{The denominator weight of the sign projector is nonzero only at the prime two.}
\label{fig:odd-dyadic}
\end{figure}

\Needspace{7\baselineskip}
\subsection{Why coefficient packets remain available}

The failure of the character formula
\[
\frac12(I\pm A)
\]
to be integral does not invalidate coefficient selectors such as even and odd residue extraction. It shows that two distinct constructions must be distinguished:
\[
\boxed{
\text{integral address selector}
\quad\neq\quad
\text{character-average spectral projector}
}
\]
when the group order is not a unit. The first can remain exact over $\Z_2$; the second can acquire a denominator weight.

\section{Smith weights and multivariable lifting}

Let $M$ be a matrix over a DVR. Smith normal form gives invertible matrices $U,V$ such that
\[
UMV
=
\diag(\pi^{s_1}u_1,\ldots,\pi^{s_r}u_r,0,\ldots,0),
\]
where each $u_i$ is a unit and
\[
0\le s_1\le\cdots\le s_r.
\]
The exponents $s_i$ are directional lifting weights.

\begin{theorem}[Concrete Smith-weighted congruence solve]\label{thm:smith-solve}
Let $p>1$ and consider the diagonal system over
$(\Z/(p^L)\Z)^d$
\[
\diag(p^{s_1},\ldots,p^{s_d})x=b
\pmod{p^L}.
\]
The system is solvable if and only if
\[
p^{\min(s_i,L)}\mid (b_i)_{\mathrm{val}}
\]
for every $i$, where $(b_i)_{\mathrm{val}}$ is the canonical natural
representative. Conditional on this admissibility condition, the number of
solutions is
\[
\boxed{
p^{\sum_{i=1}^d\min(s_i,L)}.
}
\]
If any coordinate is inadmissible, the solution set is empty and its
cardinality is zero.
\end{theorem}

\begin{proof}
Lean proves the one-coordinate arithmetic directly over
$\mathbb Z/p^L\mathbb Z$ for every $p>1$:
\[
p^s x=b\pmod{p^L}
\quad\Longleftrightarrow\quad
p^{\min(s,L)}\mid b_{\mathrm{val}},
\]
where $b_{\mathrm{val}}$ is the canonical natural representative, and every
admissible fibre has exactly $p^{\min(s,L)}$ elements. The coordinates are
independent, so the product equivalence gives the displayed global count. If
one coordinate fibre is empty, the global fibre is empty. Lean proves both the
global solvability equivalence and the conditional count, including
cardinality zero on the inadmissible branch.
\end{proof}

For an arbitrary DVR $\OO$ with residue-field cardinality $q$, the analogous
formula
\[
q^{\sum_i\min(s_i,L)}
\]
is a classical consequence of the corresponding one-coordinate solvability
and fibre-cardinality facts. In the Lean development this general form is
\emph{certificate-explicit}: those coordinate facts may be supplied to the
abstract product API, but they are not derived from a formalized general-DVR
structure theorem. Construction and diagonalization of Smith normal form for
an arbitrary DVR matrix also remain classical input; see
\cite{serre1979,lang2002}.

If a Smith form is rank deficient, a literal zero diagonal factor requires
the corresponding target coordinate to vanish modulo $p^L$ and then
contributes all $p^L$ possible coordinate values. In the concrete quotient it
can equivalently be represented by any weight $s_i\ge L$, since
$p^{s_i}=0$ in $\Z/(p^L)\Z$; the theorem then gives exactly that condition and
factor. The numbered theorem begins after the diagonal data have been
supplied and does not construct the Smith form.

\begin{example}
At $p=2$, depth $L=6$, and weights $(0,1,3)$, the right-hand side $(1,2,8)$ is admissible. The coordinate kernels have sizes
\[
1,2,8,
\]
so the system has $16$ solutions. Replacing $8$ by $4$ in the final coordinate violates the divisibility condition and makes the system unsolvable.
\end{example}

This theorem is the multivariable version of precision delay. A coordinate with weight $s_i$ contributes both a solvability threshold and $s_i$ units of branching until the ambient depth truncates it.

\section{Linear refinement defects}

Exact coordinate refinement does not guarantee that an independently chosen
coarse operator intertwines with a fine operator. Let
$C:X_f\to X_c$ be any linear coarsening map and let
$T_f\in\End(X_f)$ and $T_c\in\End(X_c)$. Define
\[
\boxed{
\Delta=CT_f-T_cC.
}
\]

\begin{theorem}[Power-defect identity]\label{thm:power-defect}
For every $n\ge1$,
\[
\boxed{
CT_f^n-T_c^nC
=
\sum_{k=0}^{n-1}
T_c^{n-1-k}\Delta T_f^k.
}
\]
\end{theorem}

\begin{proof}
The identity is immediate for $n=1$. For the general case, insert and subtract $T_cCT_f^{n-1}$:
\[
CT_f^n-T_c^nC
=(CT_f-T_cC)T_f^{n-1}
+T_c(CT_f^{n-1}-T_c^{n-1}C),
\]
and apply induction.
\end{proof}

\begin{corollary}[Norm bound]\label{cor:defect-norm}
For a submultiplicative operator norm,
\[
\|CT_f^n-T_c^nC\|
\le
\|\Delta\|
\sum_{k=0}^{n-1}
\|T_c\|^{n-1-k}\|T_f\|^k.
\]
If $\|T_f\|,\|T_c\|\le\rho$, then
\[
\boxed{
\|CT_f^n-T_c^nC\|
\le n\rho^{n-1}\|\Delta\|.
}
\]
\end{corollary}

\begin{theorem}[Resolvent-defect identity]\label{thm:resolvent-defect}
If $z-T_f$ and $z-T_c$ are invertible, then
\[
\boxed{
C(z-T_f)^{-1}-(z-T_c)^{-1}C
=
(z-T_c)^{-1}\Delta(z-T_f)^{-1}.
}
\]
\end{theorem}

\begin{proof}
Multiply the left-hand side on the left by $z-T_c$ and on the right by $z-T_f$:
\[
(z-T_c)C-C(z-T_f)=CT_f-T_cC=\Delta.
\]
Then invert the two exterior factors.
\end{proof}

For continuous linear maps, the same identity gives the certified norm bound
\[
\boxed{
\|C R_f-R_c C\|
\le \|R_c\|\,\|\Delta\|\,\|R_f\|,
}
\]
from a left inverse witness for the coarse scalar shift and a right inverse
witness for the fine scalar shift. These one-sided witnesses are sufficient
for the displayed identity, although they need not by themselves define
two-sided resolvents. If $\Delta=0$, Lean also proves, for every finite
polynomial $p$, the exact compatibility
\[
\boxed{C\,p(T_f)=p(T_c)\,C,}
\]
in both the algebraic and continuous-linear-map settings. Thus all powers,
polynomial actions, and the explicitly witnessed resolvents intertwine with
coarse reconstruction. This is the linear prototype of the more general
refinement-defect theory reserved for Paper III.

\begin{figure}[htbp]
\centering
\includegraphics[width=0.88\textwidth]{09_refinement_defect_bound.pdf}
\caption{A controlled one-step defect and its telescoping norm bound under repeated operator application.}
\label{fig:defect-bound}
\end{figure}

\section{Worked realizations in the common operator language}

\subsection{Signal polyphase coordinates}

Let $x=(x_n)_{n\in\Z}\in c_{00}(\Z;R)$ and let the seed be $\delta_0$
under the bilateral shifts $(S^r\delta_0)_n=\delta_{n,r}$. This makes the
algebraic span and every displayed sum finite. The $N$ packets are
\[
x_j[m]=x[Nm+j].
\]
Synthesis interleaves these sequences. A finite-impulse-response linear
time-invariant filter becomes an $N\times N$ matrix of transfer blocks.
Refining from $N$ to $M$ with $N\mid M$ splits each coarse phase into its
descendants without changing the reconstructed signal. The same formulas
extend to standard sequence completions when the synthesis series converges;
this is the usual polyphase viewpoint from multirate filter banks
\cite{vaidyanathan1993}.

\subsection{Transfer-block sparsity}

A banded operator need not remain banded after residue reordering, but its
nonzero couplings become a structured packet graph. With the convention
$(S_rx)_n=x_{n-r}$, the only active transfer is
\[
P_{N,j+r}S_rP_{N,j},
\]
with indices taken modulo $N$. More general finite stencils produce a finite set of active phase differences.

\begin{figure}[htbp]
\centering
\includegraphics[width=0.76\textwidth]{10_transfer_block_sparsity.pdf}
\caption{Packet-block norms for a finite stencil. The original operator is exactly reconstructed from a sparse transfer pattern.}
\label{fig:block-sparsity}
\end{figure}

\subsection{A short Smith-coordinate realization}\label{sec:smith-realization}

Normalize an impedance by $z=Z/Z_0$ and write
\[
\Gamma=\frac{z-1}{z+1}.
\]
A phase action on the reflection disk is
\[
\Gamma\mapsto\omega_N^k\Gamma.
\]
Pulling it back gives
\[
T_{N,k}(z)
=
\frac{(1+\omega_N^k)z+(1-\omega_N^k)}
{(1-\omega_N^k)z+(1+\omega_N^k)}.
\]
For $N=2$ and $k=1$,
\[
T_{2,1}(z)=\frac1z,
\qquad z\ne0,-1.
\label{eq:smith-inversion}
\]
Thus the nontrivial first-level phase action exchanges normalized impedance
and admittance. Calling these values ``packets'' is geometric terminology for
the sampled orbit, not an assertion that the two values form linear
idempotent ranges. Higher $2$-power refinements sample the same reflection
circle at progressively finer electrical phases; compare the standard Smith
chart conventions in \cite{pozar2012}.

\begin{figure}[htbp]
\centering
\includegraphics[width=0.88\textwidth]{11_smith_polyphase_orbit.pdf}
\caption{Nested $2$-power phase orbits on the Smith reflection disk. The full Smith-chart treatment is deferred to the signal paper.}
\label{fig:smith-orbit}
\end{figure}

The symmetric weight
\[
w_{\mathrm{Smith}}(\Gamma)
=
\frac1{|1-\Gamma|^2}
+
\frac1{|1+\Gamma|^2}
\]
indicates where small reflection-coordinate errors produce large impedance or admittance errors.

\begin{figure}[htbp]
\centering
\includegraphics[width=0.70\textwidth]{12_smith_adaptive_weight.pdf}
\caption{A symmetric refinement weight concentrated near the open- and short-circuit boundary points.}
\label{fig:smith-weight}
\end{figure}

\section{Numerical certificates}

All numerical work in this paper is finite and reproducible. It tests the
Python implementation of identities proved or certificate-checked above; it
does not replace those arguments or the Lean kernel checks. Release 1.2.1 runs
forty-eight tests, including invalid-input contracts, refinement transitivity,
transfer-block aggregation, Fourier/DFT agreement, brute-force Smith fibre
counts, and consistency of the stored JSON and CSV records.

\begin{table}[htbp]
\centering
\caption{Representative numerical and symbolic checks.}
\begin{tabularx}{\textwidth}{L{0.54\textwidth}Y}
\toprule
Check & Result \\
\midrule
Packet synthesis residuals for moduli $2,3,4,6,8,12$ & $0$ \\
Nested refinement residuals for $2\to4$, $2\to8$, $3\to6$, $3\to12$, $4\to8$, $4\to12$, $6\to12$ & $0$ \\
Dyadic--odd CRT residuals for $4\times3$, $8\times3$, $4\times5$, $8\times5$ & $0$ \\
Maximum finite-Fourier projector residual & $6.60\times10^{-15}$ \\
Weighted inverse residual modulo $2^{10}$ & $0$ \\
Weighted inverse residual modulo $3^8$ & $0$ \\
Weighted inverse residual modulo $5^7$ & $0$ \\
Difference between filtered and direct modular inverses & $0$ \\
Transfer-zeta coefficients versus direct walk enumeration through weight $30$ & exact match \\
Dyadic sign-projector resultant weight & $v_2(2)=1$ \\
Odd-prime sign-projector resultant weight at $3$ & $v_3(2)=0$ \\
Power-defect telescoping identity residual & at most $6.94\times10^{-18}$ \\
Resolvent-defect identity residual & $1.09\times10^{-17}$ \\
Smith $N=2$ phase versus $z\mapsto1/z$ & residual below $1.4\times10^{-16}$ \\
\bottomrule
\end{tabularx}
\end{table}

\begin{certificate}[Weighted inverse]
For the tested diagonal-plus-transfer matrices, the filtered inverse agrees exactly with direct modular inversion at all tested precisions:
\[
2^{10},\qquad3^8,\qquad5^7.
\]
The maximum modular residual is zero in every case.
\end{certificate}

\begin{certificate}[Transfer zeta]
The rational function
\[
\frac{-t^3}{t^9-t^7-t^6-t^5+2t^4+t^3+t^2-1}
\]
was expanded through degree $30$. Every coefficient agrees with direct weighted-walk enumeration.
\end{certificate}

\begin{certificate}[Dyadic separation]
The code verifies exactly that the two rational sign projectors are idempotent, orthogonal, and complementary, that scaling by $2$ makes both integral, and that the same projectors are already integral modulo every tested odd prime.
\end{certificate}

\section{Boundaries and next questions}

\begin{boundary}[Refinement is not truncation]
The exact refinement maps in this paper preserve the full object. A numerical scheme that discards fine packets requires an additional approximation theorem and cannot cite exact coordinate refinement as a truncation guarantee.
\end{boundary}

\begin{boundary}[Linear versus nonlinear defects]
The power and resolvent identities concern linear operators. For nonlinear maps, packet interactions require higher jets and operator trees. That theory is the subject of Paper III.
\end{boundary}

\begin{boundary}[Resultants give bounds]
A resultant supplies an integral denominator bound. It need not always be the minimal separation weight for a particular lattice, because cancellation can lower the actual denominator.
\end{boundary}

\begin{boundary}[Dyadic arithmetic]
The fact that a coefficient selector remains integral at $2$ does not imply that every spectral, character, or deformation projector does. Each construction must be audited for division by the group order, resultants, trace splittings, and quadratic refinement data.
\end{boundary}

The next paper will add nonlinear operation trees, exact packet-coordinate composition, refinement defects for nonlinear maps, holonomy, and obstruction packets.

\section{Conclusion}

The central result of Paper II is that refinement can be made exact without erasing the distinctions that appear in integral arithmetic. The same underlying object can be represented by a nested tower of packet coordinate systems,
\[
P_{\ell,\alpha}
=
\sum_{\beta\succ\alpha}P_{\ell+1,\beta},
\]
and every linear operator refines by summing its fine transfer blocks. Exact
coordinate changes obey identity and cocycle laws, while typed coordinate
operators satisfy a naturality square. Residue selectors are available over
arbitrary coefficient rings; a finite-Fourier realization over a commutative
semiring requires the stated finite orthogonality table, with the familiar
primitive-root formula recovered over a field. A composite modulus factors
into dyadic and odd packet coordinates, but their integral separation weights
need not agree.

Certified weighted transfer words provide the finite-quotient inverse formula
\[
q_L(T)^{-1}
=q_L\!\left(\sum_\gamma
(-1)^{|\gamma|}A_\gamma D^{-1}\right),
\]
with only $W(\gamma)<L$ retained; a complete-normed contraction supplies the
corresponding infinite Neumann inverse. Finite transfer graphs have an exact
coefficient algorithm and a semantic first-edge recurrence. For positive
integer weights, the determinant and adjugate of $I-A(t)$ automatically
identify that enumerated series with a normalized rational expression and
produce its eventual constant-coefficient recurrence. Minimal evaluated
B\'ezout hypotheses give the separator numerator identity; explicit
lattice-unit data turns it into an upward-closed separation bound. Certified
Smith coordinates convert lifting weights into solvability thresholds and
conditional branching counts, and concrete $\mathbb Z/p^L\mathbb Z$ fibres
realize those certificates. The elementary involution example shows exactly
why $2$ must be treated separately: its sign projectors have exact local
weight one at $2$, equal to the valuation of their computed resultant, and
weight zero at every odd prime.

The common operator language is therefore
\[
\boxed{
\text{packet}
\to
\text{descendants}
\to
\text{transfer words}
\to
\text{weights}
\to
\text{finite precision}
\to
\text{exact reconstruction}.
}
\]
This language will be retained in every later realization.

The release-level contribution is therefore both mathematical and
epistemic: each counted statement has a stable manuscript label, a precise
Lean target, and a declared boundary separating constructed mathematics from
supplied certificates and classical input. Kernel acceptance establishes
internal formal correctness of those targets. It does not by itself establish
novelty or replace independent mathematical review.

\appendix

\section{Notation table}

\begin{longtable}{L{0.28\textwidth}L{0.64\textwidth}}
\toprule
Symbol & Meaning \\
\midrule
\endhead
$P_{\ell,\alpha}$ & packet projector at refinement level $\ell$ \\
$I_\ell$ & packet address set at level $\ell$ \\
$\beta\succ\alpha$ & fine packet $\beta$ is a child of coarse packet $\alpha$ \\
$\cA_\ell,\cS_\ell$ & packet analysis and synthesis maps \\
$J_{\ell+1,\ell}$ & exact coordinate refinement map \\
$C_{\ell,\ell+1}$ & exact fine-to-coarse coordinate regrouping \\
$T^{(\ell)}_{\beta\alpha}$ & transfer block $P_{\ell,\beta}TP_{\ell,\alpha}$ \\
$P_{N,j}$ & residue or phase selector for address $j\pmod N$ \\
$w(e)$ & positive filtration weight of a transfer edge \\
$W(\gamma)$ & cumulative weight of a transfer word or graph walk \\
$A(t)$ & polynomial weighted adjacency matrix \\
$Z_{s,t}(t)$ & weighted transfer-zeta function from vertex $s$ to vertex $t$ \\
$s_{\pi,\Lambda}(E)$ & lattice-dependent separation weight of a rational projector, when a clearing exponent exists \\
$\Delta$ & one-step linear refinement defect $CT_f-T_cC$ \\
\bottomrule
\end{longtable}

\section{Algorithmic forms}\label{sec:algorithms}

\subsection{Residue analysis and synthesis}\label{alg:residue}

\begin{lstlisting}
def analyze_residues(x, N):
    require N > 0
    return [x[j::N] for j in range(N)]

def synthesize_residues(packets):
    N = len(packets)
    require packet lengths fit one N-way interleaving
    y = empty(sum(len(p) for p in packets))
    for j, packet in enumerate(packets):
        y[j::N] = packet
    return y
\end{lstlisting}

The verified contract is
$\operatorname{synthesize}(\operatorname{analyze}(x,N))=x$. The Lean
counterpart uses range-restricted packet coordinates and proves both inverse
laws, while the finite Python routine checks the concrete array realization.

\subsection{Weighted finite-precision inverse}\label{alg:weighted-inverse}

\begin{lstlisting}
Dinv = inverse(D)                       # certified on both sides
edge[e] = Dinv @ H_edge[e]
require Dinv @ H == sum(edge[e] for e in edges)
require edge[e] lies in filtration[weight[e]]
m = certified_lower_bound(weight)       # 1 <= m <= weight[e]
N = (L + m - 1) // m                    # so L <= m*N
S = 0
for n in range(N):
    for word in edge_tuples(edges, n):
        if sum(weight[e] for e in word) < L:
            S += (-1)**n * ordered_product(edge[e] for e in word)
return q_L(S @ Dinv)
\end{lstlisting}

Lean proves that the mapped pruned sum is exactly the mapped finite geometric
inverse. The lower bound $m$ gives $m|\gamma|\le W(\gamma)$, so
$L\le mN$ derives $q_L(A^N)=0$ from the filtration certificate and makes the
returned value a two-sided inverse of $q_L(D+H)$. The Lean lemma
$L\le m\lceil L/m\rceil$ justifies the displayed executable cutoff for
$m>0$. Taking $m=1$ and $N=L$ gives
the unconditional positive-weight cutoff. The Python implementation checks the
DVR matrix specialization against direct modular inversion.

\subsection{Transfer-zeta walk construction}\label{alg:path-zeta}

\begin{lstlisting}
def coefficient(G, start, target, n):
    total = 0
    for length in range(n + 1):
        for word in edge_tuples(G, length):
            if is_walk(word, start, target) and weight(word) == n:
                total += 1
    return total

A[i,j] = sum(t**weight(e) for distinct edges e from i to j)
T = I - A
Q = det(T)
H = adjugate(T)
P = H[start, target]
assert Q.coeff(0) == 1
Z = coefficient_series(G, start, target)
prove Q * Z == P
\end{lstlisting}

Positive weights make the coefficient loop finite because every contributing
word has length at most $n$. Lean proves the enumeration correct, derives its
first-edge recurrence, proves $(I-A)Z=I$, constructs the determinant/adjugate
inverse, and identifies the two matrices by inverse uniqueness. Thus the
normalized $QZ=P$ identity and its induced recurrence follow from finite graph
data. The Python companion repeats the exact symbolic construction as an
executable cross-check; it is not an oracle used by the Lean proof.

\section{Proof dependency map}

\begin{figure}[htbp]
\centering
\includegraphics[width=0.96\textwidth]{13_theorem_dependency.pdf}
\caption{The dependency structure of Paper II. Coordinate coherence, residue/CRT, weighted inversion, automatic finite-graph transfer rationality, and integral separation are distinct branches with explicit inputs.}
\label{fig:dependency}
\end{figure}

The proof order is:
\[
\begin{aligned}
&\text{complete packet axioms}
  \Rightarrow\text{finite-packet uniqueness and coordinate coherence}\\
&\qquad\Rightarrow\text{coordinate naturality and refined transfer blocks},\\
&\text{residue selectors and a finite Fourier table}
  \Rightarrow\text{Fourier formulas},\\
&\text{coprime moduli}\Rightarrow\text{dyadic--odd formulas},\\
&\text{quotient, filtration, and edge-weight certificates}
  \Rightarrow\text{explicit pruned inverse},\\
&\text{finite positive-weight graph}
  \Rightarrow\text{semantic enumeration and first-edge recurrence}\\
&\qquad\Rightarrow\det/\operatorname{adj}\text{ rationality and recurrence},\\
&\text{minimal evaluated B\'ezout data and inverse-unit lattice action}\\
&\qquad\Rightarrow\text{separation bound, upward closure, and precision delay},\\
&\text{concrete }\mathbb Z/p^L\mathbb Z\text{ coordinate fibres}
  \Rightarrow\text{finite Smith product count},\\
&\text{one-step typed defect}
  \Rightarrow\text{powers, polynomial actions, and witnessed resolvents}.
\end{aligned}
\]

The transfer-zeta walk branch is logically independent of operator inversion. Its
rationality conclusion now follows from finite graph data, rather than from a
supplied $QZ=P$ premise. Likewise, CRT factorization follows directly from
residue selectors and coprimality; it does not require the abstract refinement
theorem, although it supplies a concrete refinement family. General Smith
diagonalization and construction of polynomial B\'ezout/resultant data remain
inputs upstream of the formalized post-certificate branches.

\section{Reproducibility layout}

\begin{lstlisting}
Projector_Resolved_Seed_Calculus_II_v1.2.1/
  paper/
    Projector_Resolved_Seed_Calculus_II.{tex,pdf}
    figures/                         # 13 PDF/PNG pairs
  formalization/lean/
    KomorebiVerification/PaperII/    # modules + KernelAudit
    lean-toolchain, lake-manifest.json
    LEAN_THEOREM_MAP_v1.1.1.md
  code/python/
    seed_refinement/, tests/, dist/
    pyproject.toml, requirements.{in,lock}
  results/                           # JSON, CSV, run logs
  provenance/                        # claim map and audits
  README.md, MANIFEST.md, RELEASE.json, CHANGELOG.md
  CITATION.cff, CITATION.bib, LICENSE_STATUS.md
  Run_All.{sh,ps1}, SHA256SUMS.txt
\end{lstlisting}

The Lean toolchain and complete dependency graph are pinned. The source audit
rejects placeholders and project-specific axioms, the Python run uses an exact
lock file, and the root checksum manifest covers every release file other than
the manifest itself. The ongoing series checklist remains separate.

\begin{thebibliography}{99}

\bibitem{komorebi2026projectorI}
Komorebi P.,
``Projector-Resolved Seed Calculus I: Typed Seeds, Operator Grammars, and Exact
Reconstruction,'' version 1.2.1, manuscript and verification release, 2026.

\bibitem{demoura2015lean}
L. de Moura, S. Kong, J. Avigad, F. van Doorn, and C. Roux,
``The Lean theorem prover (system description),'' in
\emph{Automated Deduction---CADE-25}, Lecture Notes in Computer Science 9195,
pp. 378--388, Springer, 2015.
\url{https://doi.org/10.1007/978-3-319-21401-6_26}.

\bibitem{mathlib2020}
The Mathlib Community,
``The Lean mathematical library,'' in
\emph{Proceedings of the 9th ACM SIGPLAN International Conference on
Certified Programs and Proofs}, pp. 367--381, 2020.
\url{https://doi.org/10.1145/3372885.3373824}.

\bibitem{serre1979}
J.-P. Serre, \emph{Local Fields}, Graduate Texts in Mathematics 67,
Springer, 1979.
\url{https://doi.org/10.1007/978-1-4757-5673-9}.

\bibitem{lang2002}
S. Lang, \emph{Algebra}, 3rd ed., Graduate Texts in Mathematics 211,
Springer, 2002.
\url{https://doi.org/10.1007/978-1-4613-0041-0}.

\bibitem{stanley2012}
R. P. Stanley, \emph{Enumerative Combinatorics}, vol. 1, 2nd ed.,
Cambridge University Press, 2012.
\url{https://doi.org/10.1017/CBO9781139058520}.

\bibitem{vaidyanathan1993}
P. P. Vaidyanathan, \emph{Multirate Systems and Filter Banks},
Prentice Hall, 1993. ISBN 978-0-13-605718-5.

\bibitem{pozar2012}
D. M. Pozar, \emph{Microwave Engineering}, 4th ed., Wiley, 2012.
ISBN 978-0-470-63155-3.

\end{thebibliography}

\end{document}
